\documentclass[10pt, letterpaper]{article}

% Packages:
\usepackage[
    ignoreheadfoot,
    top=0.4 cm,
    bottom=0.4 cm,
    left=0.8 cm,
    right=0.8 cm,
    footskip=0.2 cm,
]{geometry}
\usepackage{titlesec}
\usepackage{tabularx}
\usepackage{array}
\usepackage[dvipsnames]{xcolor}
\definecolor{primaryColor}{RGB}{0, 0, 0}
\usepackage{enumitem}
\usepackage{amsmath}
\usepackage[
    pdftitle={Tanmay Bhuskute's CV},
    pdfauthor={Tanmay Bhuskute},
    pdfcreator={LaTeX with RenderCV},
    colorlinks=true,
    urlcolor=primaryColor
]{hyperref}
\usepackage[pscoord]{eso-pic}
\usepackage{calc}
\usepackage{bookmark}
\usepackage{lastpage}
\usepackage{changepage}
\usepackage{paracol}
\usepackage{ifthen}
\usepackage{needspace}
\usepackage{iftex}

\ifPDFTeX
    \input{glyphtounicode}
    \pdfgentounicode=1
    \usepackage[T1]{fontenc}
    \usepackage[utf8]{inputenc}
    \usepackage{lmodern}
\fi

\usepackage{XCharter}

% Some settings:
\raggedright
\AtBeginEnvironment{adjustwidth}{\partopsep0pt}
\pagestyle{empty}
\setcounter{secnumdepth}{0}
\setlength{\parindent}{0pt}
\setlength{\topskip}{0pt}
\setlength{\columnsep}{0.15cm}
\pagenumbering{gobble}

% ── Section title ──
\titleformat{\section}{\needspace{2\baselineskip}\bfseries\large}{}{0pt}{}[\vspace{1pt}\titlerule]
\titlespacing{\section}{-1pt}{0.15 cm}{0.08 cm}

\renewcommand\labelitemi{$\vcenter{\hbox{\small$\bullet$}}$}

% ── Bullet lists ──
\newenvironment{highlights}{
    \begin{itemize}[
        topsep=0.10 cm,
        parsep=0.06 cm,
        partopsep=0pt,
        itemsep=0.06 cm,
        leftmargin=12pt
    ]
}{
    \end{itemize}
}

\newenvironment{highlightsforbulletentries}{
    \begin{itemize}[
        topsep=0.10 cm,
        parsep=0.06 cm,
        partopsep=0pt,
        itemsep=0.04 cm,
        leftmargin=12pt
    ]
}{
    \end{itemize}
}

\newenvironment{onecolentry}{
    \begin{adjustwidth}{0 cm + 0.00001 cm}{0 cm + 0.00001 cm}
}{
    \end{adjustwidth}
}

\newenvironment{twocolentry}[2][]{
    \onecolentry
    \def\secondColumn{#2}
    \setcolumnwidth{\fill, 4.5 cm}
    \begin{paracol}{2}
}{
    \switchcolumn \raggedleft \secondColumn
    \end{paracol}
    \endonecolentry
}

\newenvironment{threecolentry}[3][]{
    \onecolentry
    \def\thirdColumn{#3}
    \setcolumnwidth{, \fill, 4.5 cm}
    \begin{paracol}{3}
    {\raggedright #2} \switchcolumn
}{
    \switchcolumn \raggedleft \thirdColumn
    \end{paracol}
    \endonecolentry
}

\newenvironment{header}{
    \setlength{\topsep}{0pt}\par\kern\topsep\centering\linespread{0.6}
}{
    \par\kern\topsep
}

\let\hrefWithoutArrow\href

\begin{document}
    \newcommand{\AND}{\unskip
        \cleaders\copy\ANDbox\hskip\wd\ANDbox
        \ignorespaces
    }
    \newsavebox\ANDbox
    \sbox\ANDbox{$|$}

    \begin{header}
        \fontsize{16 pt}{24 pt}\selectfont Tanmay Bhuskute

        \vspace{3 pt}

        \normalsize
        \mbox{ +1 (732) 858-3027 }%
        \kern 1.0 pt%
        \AND%
        \kern 1.0 pt%
        \mbox{ tanmay.v.bhuskute@gmail.com }%
        \kern 1.0 pt%
        \AND%
        \kern 1.0 pt%
        \mbox{ github.com/Tanmayb05 }%
        \kern 1.0 pt%
        \AND%
        \kern 1.0 pt%
        \mbox{ linkedin.com/in/tanmay-bhuskute }%
        \kern 1.0 pt%
        \AND%
        \kern 1.0 pt%
        \mbox{ Cincinnati, OH }%
    \end{header}

    \vspace{2 pt}

%═══════════════════════════════════════
\section{SUMMARY}
%═══════════════════════════════════════
\begin{onecolentry}
Gen AI and Python engineer building production LLM applications: agentic orchestration, tool calling, structured output extraction with JSON schema and Pydantic validation, RAG and vector retrieval, and FastAPI microservices on AWS, Azure, and GCP. Strong in async Python, SQL, pytest, Docker, and CI/CD.
\end{onecolentry}

%═══════════════════════════════════════
\section{EDUCATION}
%═══════════════════════════════════════
\begin{twocolentry}{
    Aug 2024 - Dec 2026
}
    \textbf{Master of Science in Computer Science} \\
    University of Cincinnati
\end{twocolentry}
\begin{onecolentry}
Relevant Coursework: Deep Learning, Intelligent Data Analysis, Database Theory.
\end{onecolentry}
\vspace{0.15 cm}

\begin{twocolentry}{
    Aug 2018 - Jul 2022
}
    \textbf{Bachelor of Engineering in Computer Engineering} \\
    AISSMS College of Engineering
\end{twocolentry}
\begin{onecolentry}
Relevant Coursework: Machine Learning, Data Mining, Object-Oriented Programming.
\end{onecolentry}
%═══════════════════════════════════════
\section{TECHNICAL SKILLS}
%═══════════════════════════════════════
\begin{onecolentry}
\textbf{Programming:} Python, TypeScript, JavaScript, Java, SQL, Go, C/C++, Bash.
\end{onecolentry}
\vspace{0.06 cm}
\begin{onecolentry}
\textbf{Gen AI \& Agentic:} LLM Agents, Function Calling, Tool Use, MCP, Prompt Engineering, LangChain, RAG, Embeddings, Chunking, Vector Databases, Structured Output, JSON Schema, Pydantic, Guardrails, Evaluation Harnesses.
\end{onecolentry}
\vspace{0.06 cm}
\begin{onecolentry}
\textbf{Model APIs \& ML:} Google Gemini, Vertex AI, GPT-4o, Llama, Mistral, Ollama, PyTorch, scikit-learn, pandas, NumPy.
\end{onecolentry}
\vspace{0.06 cm}
\begin{onecolentry}
\textbf{Backend \& Data:} FastAPI, Flask, Django, Node.js, Express, REST APIs, Microservices, Async Python, PostgreSQL, MySQL, Redis, PySpark, Delta Lake.
\end{onecolentry}
\vspace{0.06 cm}
\begin{onecolentry}
\textbf{Cloud \& Tooling:} Azure, AWS, GCP, Docker, Kubernetes, GitLab CI/CD, Git, Terraform, Linux, pytest, SonarQube.
\end{onecolentry}

%═══════════════════════════════════════
\section{WORK EXPERIENCE}
%═══════════════════════════════════════
\begin{twocolentry}{Feb 2026 -- Present}
\textbf{Graduate Researcher, Gen AI} \,|\, University of Cincinnati
\end{twocolentry}
\begin{highlights}
  \item Built two-stage \textbf{agentic LLM pipeline} in Python over \textbf{Gemini 2.5 Pro and Vertex AI}, extracting a structured task memory once, then conditioning a multi-step closed-loop automation agent on it instead of re-running per-step model inference.
  \item Delivered \textbf{25.7x latency reduction} (39.4s vs 1012.1s avg per run) against the published \textbf{ViBR} GPT-4o baseline while preserving task-completion fidelity, measured across \textbf{53 runs on 20+ apps}.
  \item Built config-driven \textbf{evaluation harness} and \textbf{provider-agnostic model abstraction} (Gemini, GPT-4o, Llama, Qwen) making prompt and model changes measurable via reproducible regression runs.
\end{highlights}

\vspace{0.15 cm}

\begin{twocolentry}{Jul 2022 -- Jun 2024}
\textbf{Software Development Engineer} \,|\, Siemens Industry Software
\end{twocolentry}
\begin{highlights}
  \item Wrote production-grade \textbf{Python (Boto3)} automation replacing legacy Ruby services, cutting deployment time \textbf{40\%} across AWS environments and shipping through peer code review and CI/CD.
  \item Executed \textbf{44 production upgrades} in 4 months, debugging live services across Linux, EC2, and databases; designed an \textbf{Ansible Blue-Green} framework that cut upgrade downtime \textbf{50\%}.
  \item Drove engineering quality as \textbf{Quality Coach and SonarQube Specialist}, raising automated test coverage from \textbf{35\% to 78\%} via coding standards and CI quality gates; mentored \textbf{5 engineers} and migrated \textbf{40+ environments} to Kubernetes with ArgoCD, Helm, and GitLab CI/CD.
\end{highlights}

\vspace{0.15 cm}

\begin{twocolentry}{Mar 2022 -- Jun 2022}
\textbf{Software Developer Intern} \,|\, Siemens Industry Software
\end{twocolentry}
\begin{highlights}
  \item Automated AWS provisioning with \textbf{Terraform}, cutting deployment time \textbf{30\%}; built a \textbf{Go and Terratest} validation suite catching \textbf{95\%} of misconfigurations before production.
\end{highlights}

%═══════════════════════════════════════
\section{PROJECTS AND RESEARCH}
%═══════════════════════════════════════
\begin{twocolentry}{May 2026}
\textbf{PolicyPilot -- Policy Intelligence Portal} [Python, Gemini, MCP, RAG, TS, Node]
\end{twocolentry}
\begin{highlights}
  \item Architected evidence-grounded \textbf{RAG pipeline} over \textbf{Gemini 2.5-flash} with chunking, retrieval, and citation-backed answers, eliminating hallucination in high-stakes eligibility decisions.
  \item Designed \textbf{17-tool MCP server} exposing structured \textbf{LLM tool-calling} with \textbf{schema validation} at every data boundary, plus \textbf{confidence-scoring and disambiguation} logic (evidence count, negation detection, subjectivity flags) routing uncertain cases to human review.
\end{highlights}

\vspace{0.12 cm}

\begin{twocolentry}{Jul 2026}
\textbf{NYC Mobility Forecasting} [Python, PySpark, Delta Lake, FastAPI, PostgreSQL]
\end{twocolentry}
\begin{highlights}
  \item Built \textbf{FastAPI} prediction microservice on PostgreSQL and Redis caching, containerized with Docker, over a \textbf{2.3M-row} Spark dataset; two-stage model improved MAPE \textbf{18.77\% to 16.52\%}.
\end{highlights}

%═══════════════════════════════════════
\section{EXTRACURRICULAR ACTIVITIES}
%═══════════════════════════════════════
\begin{onecolentry}
\begin{highlights}
  \item \textbf{Finalist, RevolutionUC 2025 Hackathon} (300 participants, 31 schools) \,|\, \textbf{SAFe Practitioner 5.1} (Mar 2023).
\end{highlights}
\end{onecolentry}

\end{document}
